\documentclass[12pt,a4paper]{report}

% --- Pacotes Básicos ---
\usepackage[utf8]{inputenc}
\usepackage[T1]{fontenc}
\usepackage[brazil]{babel}
\usepackage{graphicx}
\usepackage{geometry}
\geometry{a4paper, left=3cm, top=3cm, right=2cm, bottom=2cm}
\usepackage{setspace}
\onehalfspacing
\usepackage{indentfirst}
\setlength{\parindent}{1.25cm}
\usepackage{titlesec}
\usepackage{booktabs}
\usepackage{tabularx}

% --- Configuração de Títulos ---
\titleformat{\chapter}[hang]{\normalfont\huge\bfseries}{\thechapter}{1em}{}
\titlespacing*{\chapter}{0pt}{-20pt}{20pt}

% --- Dados do Projeto (Preencher aqui) ---
\newcommand{\projetoNome}{[NOME DO PROJETO]}
\newcommand{\versao}{1.0}
\newcommand{\data}{\today}

\begin{document}

% --- Capa Simples ---
\begin{titlepage}
    \begin{center}
        \large UNIVERSIDADE CATÓLICA DE BRASÍLIA \\
        \large CURSO DE CIÊNCIA DA COMPUTAÇÃO \\
        \vspace{4cm}
        \huge \textbf{Documento de Visão} \\
        \vspace{1cm}
        \Large \projetoNome \\
        \vspace{4cm}
        \large Versão: \versao \\
        \vfill
        \large Brasília, DF \\
        \large \data
    \end{center}
\end{titlepage}

\tableofcontents
\newpage

\chapter{Introdução}
[Apresente aqui o propósito deste documento, qual o grupo e seus respectivos membros. Explique também se este é o projeto de construção de um novo software, qual a versão que será construída e se há um cliente externo, se é um software multicliente para mercado ou sem fins lucrativos. Destaque a motivação e importância do desenvolvimento da solução proposta.]

\chapter{Contexto de Negócio}
[Descreva detalhadamente o cenário atual do negócio ou domínio de aplicação, sem abordar a solução proposta. Identifique e analise as principais deficiências e lacunas enfrentadas no contexto, destacando problemas, ineficiências ou necessidades não atendidas. Apresente as ferramentas ou métodos atualmente utilizados, indicando suas limitações.]

\section{Público e Mercado}
[Descreva os principais usuários e/ou clientes impactados por esse contexto, bem como o mercado-alvo. Inclua indicadores quantitativos ou qualitativos que evidenciem a relevância e urgência do problema.]

\section{Benchmarking}
[Realize um benchmarking com soluções ou softwares existentes que atendam esse mercado ou domínio, comparando suas funcionalidades, pontos fortes e fracos.]

\chapter{Posicionamento}

\section{Declaração do Problema}
\begin{center}
\begin{tabularx}{\textwidth}{|l|X|}
    \hline
    \textbf{O problema de} & [descreva o problema] \\ \hline
    \textbf{afeta} & [as partes interessadas afetadas pelo problema] \\ \hline
    \textbf{cujo impacto é} & [qual é o impacto do problema?] \\ \hline
    \textbf{uma solução de sucesso deveria} & [liste alguns benefícios chave de uma solução de sucesso] \\ \hline
\end{tabularx}
\end{center}

\section{Declaração da Visão do Software}
\begin{center}
\begin{tabularx}{\textwidth}{|l|X|}
    \hline
    \textbf{Para} & [clientes e usuários alvo] \\ \hline
    \textbf{Que} & [declaração da necessidade ou oportunidade] \\ \hline
    \textbf{O} & \projetoNome \\ \hline
    \textbf{É um} & [categoria do produto] \\ \hline
    \textbf{Que} & [objetivo principal do software; do ponto de vista do usuário] \\ \hline
    \textbf{Diferente de} & [alternativa do mercado, principal concorrente ou solução atual] \\ \hline
    \textbf{Nosso produto} & [principais diferenciais] \\ \hline
\end{tabularx}
\end{center}

\chapter{Descrição das Partes Interessadas}
\begin{center}
\begin{tabularx}{\textwidth}{|l|X|X|}
    \hline
    \textbf{Nome} & \textbf{Descrição} & \textbf{Responsabilidades} \\ \hline
    [Nome da Parte] & [Breve descrição] & [Resumo das responsabilidades] \\ \hline
\end{tabularx}
\end{center}

\chapter{Visão Geral do Produto}

\section{Necessidades e Funcionalidades}
\begin{center}
\begin{tabularx}{\textwidth}{|X|X|l|l|}
    \hline
    \textbf{Necessidade} & \textbf{Funcionalidade} & \textbf{Prioridade} & \textbf{Responsável} \\ \hline
    & & & \\ \hline
\end{tabularx}
\end{center}

\section{Requisitos Não Funcionais Preliminares}
\begin{center}
\begin{tabularx}{\textwidth}{|X|l|}
    \hline
    \textbf{Requisito não funcional} & \textbf{Prioridade} \\ \hline
    & \\ \hline
\end{tabularx}
\end{center}

\end{document}
